\documentclass[11pt,a4paper]{article}

\usepackage[utf8]{inputenc}
\usepackage[T1]{fontenc}
\usepackage{lmodern}
\usepackage[margin=1in]{geometry}
\usepackage{booktabs}
\usepackage{array}
\usepackage{graphicx}
\usepackage{hyperref}
\usepackage{xcolor}
\usepackage{amsmath}
\usepackage{natbib}

\definecolor{linkblue}{RGB}{31, 119, 180}
\hypersetup{colorlinks=true,linkcolor=linkblue,citecolor=linkblue,urlcolor=linkblue}

\setlength{\parskip}{0.5em}
\setlength{\parindent}{0em}

\begin{document}

\begin{center}
{\LARGE\bfseries Beyond Recall: Behavioral Specification as the\\
Missing Primitive for AI Personalization}

\vspace{1em}
{\large Aarik Gulaya} \\
{\normalsize Base Layer} \\[0.5em]
{\small April 2026 \quad $\cdot$ \quad Preprint \quad $\cdot$ \quad Apache 2.0}\\[0.3em]
{\small \href{https://github.com/agulaya24/base-layer}{github.com/agulaya24/base-layer} \quad $\cdot$ \quad
\href{https://github.com/agulaya24/memory-study-repo}{github.com/agulaya24/memory-study-repo}}

\vspace{0.3em}
{\small\textit{Disclosure: The author is the founder of Base Layer, which provides an open-source
implementation of the pipeline described in this paper.
All data, code, and evaluation artifacts are released under Apache 2.0.}}
\end{center}

\vspace{1em}
\hrule
\vspace{1em}

\begin{abstract}
In order to have aligned AI memory and actions---where alignment means the AI's behavior
accords with how a specific person reasons and decides---AI requires a representation of how
its user thinks. Current memory systems store what someone said. Preference models store what
someone likes. Neither captures how someone reasons.

We introduce \textbf{Behavioral Specifications}: a compressed document ($\sim$5K tokens)
encoding how a person thinks, decides, and reasons. It is generated from source text by
extracting behavioral facts using 47 constrained predicates, authoring three interpretive
layers (anchors, core, predictions), and composing them into a unified specification.

We tested this across 14 subjects from 11 cultures, 6 response models, 4 state-of-the-art
commercial memory systems (Mem0, Letta, Supermemory, Zep), 3 providers, and 7 calibrated
LLM judges. The behavioral specification significantly improves representational accuracy
for subjects with low pretraining representation (12 of 14 subjects, Wilcoxon signed-rank
$p=0.0015$), with improvements ranging from $+6\%$ to $+136\%$ over baseline.

\vspace{0.5em}
\noindent Key findings:

\begin{enumerate}
\item Adding a behavioral specification to 4 SOTA memory systems increases representational
accuracy by $+12\%$ to $+66\%$ (Letta $+45\%$, Mem0 $+22\%$, Supermemory $+12\%$, Zep $+66\%$).

\item A 5,250-word ($\sim$8.6K token) spec outperforms 25,000 words ($\sim$33K tokens) of raw
source text by 29\% normalized.

\item For low-representation subjects (baseline below 2.4/5.0), improvement ranges from
$+6\%$ to $+136\%$. For well-represented subjects, the specification is unnecessary or
slightly harmful.

\item Four SOTA systems have zero top-1 retrieval overlap 68\% of the time given the same
facts and the same question. The specification still improves every system.

\item A wrong behavioral specification scores near baseline across all 14 subjects. The
correct content drives the improvement, not the presence of a framework.
\end{enumerate}

\noindent Seven judges confirm the same condition rankings (pairwise Spearman
$\rho = 0.89$--$0.98$). Krippendorff's $\alpha = 0.723$.
\end{abstract}

\vspace{1em}
\hrule
\vspace{1em}

\section{1. Introduction}

In order to have aligned AI memory and actions---where alignment means the AI's behavior
accords with how a specific person reasons and decides---AI requires a representation of how
its user thinks. Current memory systems store what someone said. Preference models store what
someone likes. Neither captures how someone reasons.

We call this a \textbf{Behavioral Specification}: a compressed document ($\sim$5K tokens)
that encodes a person's decision patterns, values under conflict, risk tolerance, learning
style, and reasoning structure. It is not a preference model, not a fact store, and not a
persona. It is a model of how someone reasons about their experience, generated from source
text through a fully automated pipeline (Section 3.3).

The definitive memory benchmark, LongMemEval~\citep{he2025longmemeval} (ICLR 2025), tests
recall: ``What did the user say about X in conversation 47?'' Four state-of-the-art commercial
memory systems (Mem0, Letta, Supermemory, Zep) all score 85\%+ on these retrieval tasks.
But recall and preference storage share the same limitation: they capture what was expressed,
not the reasoning behind it.

To illustrate the gap: when we gave four SOTA memory systems the same 462 facts and the same
question, they retrieved completely different top-1 facts 68\% of the time (39\% at top-3,
22\% at top-5). They all pass recall benchmarks, but cannot agree on what is relevant.

\subsection{1.1 What We Mean by Alignment}

In this study, ``alignment'' refers to \textit{behavioral alignment}: making an AI's actions
accord with a specific person's reasoning patterns, values, and decision-making. We ask:
can the AI act the way you would act, given how you think?

\subsection{1.2 What the Specification Captures}

The Behavioral Specification identifies \textit{interpretive patterns}: durable structures
of reasoning that show up over time. A memory system that stores ``his father was violent
(biographical, from Hamerton's autobiography)'' has captured a fact. A specification that
captures ``this person evaluates authority figures on two simultaneous ledgers, virtue and
failure, refusing to collapse them into a single verdict (behavioral, derived from the same
source)'' has captured the interpretive pattern that gives that fact its significance. The
pattern was never stated. It was derived from dozens of facts. This is what makes it a
specification rather than a snapshot.

% --- Example Table ---
\section{4. Results (Excerpt)}

\subsection{4.4 The Global Gradient: N=14 Across 11 Cultures}

\begin{table}[h]
\centering
\caption{Global gradient: specification impact across 14 subjects. Scores are means
across 6 judges. Effect: percentage improvement of Spec over Baseline.}
\label{tab:gradient}
\small
\begin{tabular}{llrrrrrrr}
\toprule
\textbf{Subject} & \textbf{Culture} & \textbf{Baseline} & \textbf{Spec} &
\textbf{Wrong} & \textbf{Facts} & \textbf{F+Spec} & \textbf{Effect} & \textbf{Abs.} \\
\midrule
Hamerton        & British          & 1.25 & 2.94 & 1.79 & 2.55 & 3.08 & \textbf{+136\%} & +1.69 \\
Sunity Devee    & Indian           & 1.03 & 2.27 & 1.29 & 2.46 & 2.41 & \textbf{+121\%} & +1.24 \\
Georg Ebers     & German           & 1.04 & 1.80 & 1.50 & 2.21 & 2.34 & \textbf{+73\%}  & +0.76 \\
Fukuzawa        & Japanese         & 1.80 & 2.56 & 2.11 & 2.89 & 2.99 & \textbf{+42\%}  & +0.75 \\
Mary Seacole    & Caribbean        & 1.77 & 2.48 & 1.43 & 2.63 & 2.60 & \textbf{+40\%}  & +0.71 \\
Bernal Diaz     & Latin American   & 1.83 & 2.53 & 2.13 & 2.71 & 2.78 & \textbf{+38\%}  & +0.70 \\
Keckley         & Black American   & 1.92 & 2.64 & 1.50 & 2.57 & 2.62 & \textbf{+38\%}  & +0.73 \\
Yung Wing       & Chinese          & 1.88 & 2.22 & 2.20 & 2.13 & 2.40 & \textbf{+18\%}  & +0.34 \\
Rousseau        & French           & 2.44 & 2.81 & 1.91 & 2.32 & 2.53 & \textbf{+15\%}  & +0.37 \\
Babur           & Central Asian    & 2.03 & 2.28 & 1.23 & 2.37 & 2.39 & \textbf{+12\%}  & +0.25 \\
Cellini         & Italian          & 2.56 & 2.72 & 1.94 & 2.61 & 2.80 & \textbf{+6\%}   & +0.16 \\
Augustine       & North African    & 2.80 & 2.97 & 2.54 & 3.09 & 3.22 & \textbf{+6\%}   & +0.17 \\
Equiano         & West African     & 2.93 & 2.70 & 2.18 & 2.63 & 2.65 & $-8\%$          & $-0.24$ \\
Zitkala-Sa      & Native American  & 2.34 & 2.03 & 1.66 & 2.10 & 2.02 & $-13\%$         & $-0.31$ \\
\bottomrule
\end{tabular}

\vspace{0.3em}
\noindent\small\textit{Wilcoxon signed-rank: $p=0.0015$, $N=14$.
Krippendorff's $\alpha=0.723$. Effect: relative improvement over baseline raw score.
Abs.: absolute score difference.}
\end{table}

\section{8. Conclusion}

The problem is not recall. Four memory systems score 85\%+ on the same benchmark. Given the
same facts and the same question, they retrieve completely different facts 68\% of the time.
They have solved storage. They have not solved understanding.

The problem is interpretation. A fact does not carry its own significance. The lens through
which it is evaluated is what gives it personal significance. No memory system captures that
lens. The behavioral specification does.

The specification provides for private individuals what pretraining provides for public
figures. It does so with full traceability. This is what separates personalization from
profiling.

Memory systems store what was said. Preference models store what was liked. The behavioral
specification captures how someone makes sense of what happened to them. That is the
primitive that has been missing.

\section*{Acknowledgments}

This research was self-funded. No external funding, grants, or institutional support was
received. The author thanks the research community whose work made this study possible,
and his wife Bavani, his mother, Walnut, and Wasabi for their unwavering support, which
allowed this work to come into existence.

\end{document}
